% Source: https://artofproblemsolving.com/community/c3h1558995_infinity
% Stable identity: aops-topic-1558995-q01
Let $N=(n+1)!$ and consider the $n$ consecutive integers
\[
N+2,\ N+3,\ \dots,\ N+(n+1).
\]
For each $k$ with $2\le k\le n+1$, the number $k$ divides $N=(n+1)!$, so $k$ divides $N+k$. Also $N+k>k$, so $N+k$ has the divisor $k$ with $1<k<N+k$ and is not prime. Hence all $n$ of these numbers are composite. $\blacksquare$

For example, with $n=4$: $5!=120$ gives $122,123,124,125$, divisible by $2,3,4,5$ respectively. (The gaps found this way are far from the first ones; for instance $24,25,26,27,28$ are already five consecutive composites. Note also that "infinity" is not an integer, so it is neither prime nor composite.)
